% Guia do professor — Capítulo 22: Óptica Atmosférica
\section{Capítulo 22 --- Óptica Atmosférica}

\subsection*{Visão geral e ritmo}

O capítulo-sobremesa do livro: Snell, dispersão, dipolos e refração
aplicados ao céu. Para físicos é território natal — e a mensagem
docente é que cada fenômeno luminoso é um \emph{diagnóstico} (halo $=$
cirrus; arco $=$ gotas grandes; miragem $=$ perfil de $T$).
\textbf{Ritmo: 2 aulas de 2\,h (rende bem como fecho leve do
semestre).}

\begin{itemize}
  \item \textbf{Aula 1} — arco-íris completo (cáustica no quadro,
        Monte Carlo no Caso~1), halos e a família dos cristais
        (Caso~2).
  \item \textbf{Aula 2} — Rayleigh e as cores do céu (Caso~3),
        polarização, refração e miragens (Caso~4), flash verde, e a
        síntese do curso.
\end{itemize}

\subsection*{Objetivos de aprendizagem}

\begin{enumerate}
  \item deduzir o ângulo do arco-íris pela cáustica do desvio mínimo;
  \item explicar o halo de 22° pelo desvio mínimo do prisma de gelo;
  \item usar Rayleigh ($\lambda^{-4}$) $+$ Beer para céu azul e poente
        vermelho;
  \item calcular a curvatura de raios em gradientes de $n$ e
        classificar miragens;
  \item ler fenômenos ópticos como diagnósticos de nuvens e perfis.
\end{enumerate}

\subsection*{Pré-requisitos a reativar}

Snell e dispersão (óptica básica); Beer e massa de ar (Cap.~2); gotas
e cristais e seus tamanhos (Caps.~6--7); inversões (Cap.~5); lei dos
gases para $n(\rho)$ (Cap.~1); dipolo irradiante (eletromagnetismo).

\subsection*{Armadilhas conceituais frequentes}

\begin{itemize}
  \item \textbf{``O arco está num lugar''}: é um cone de 42° em torno
        do antissolar de CADA observador — não há pé do arco-íris.
  \item \textbf{Halo confundido com coroa}: halo $=$ refração em gelo
        (22°, borda vermelha interna); coroa $=$ difração em gotículas
        (poucos graus, vermelho externo).
  \item \textbf{``Céu azul porque reflete o mar''}: o dipolo de
        Rayleigh resolve em uma linha (T3) — e a polarização (T5) é a
        prova experimental barata (óculos polarizados!).
  \item \textbf{Miragem como ``ilusão''}: os raios chegam de verdade
        por caminhos curvos — a imagem é real, a interpretação do
        cérebro é que assume linha reta.
  \item \textbf{Flash verde como lenda}: é dispersão $+$ refração
        padrão; raro porque exige horizonte limpo e baixo aerossol.
\end{itemize}

\subsection*{Demonstração de baixo custo}

Laser $+$ aquário com gelatina/açúcar em gradiente: o raio curva à
vista (miragem de bancada). Óculos polarizados girando contra o céu a
90° do Sol. Esfera de vidro/gota pendurada ao sol para o arco de
bancada. CD como rede de difração para o espectro do céu.

\subsection*{Problemas sugeridos do livro (exercícios do Cap.~22)}

Aquecimento: ângulos de arco e halo com $n$ dados; massa de ar do
poente. Aprofundamento: curvatura de raios com gradientes dados;
polarização. Síntese: ``você vê um halo pela manhã — que nuvem é, o
que vem nas próximas 24\,h e por quê?'' (amarra Caps.~6, 12 e 22).
\emph{Confira a numeração na sua cópia (v1.02b).}

\subsection*{Uso do notebook \texttt{cap22\_optica.ipynb}}

\begin{center}
\begin{tabular}{lll}
\hline
Caso & Conteúdo & Momento sugerido \\
\hline
1 & arco-íris: traçado de raios $+$ Monte Carlo & Aula 1 \\
2 & halo 22°: prisma e histograma de cristais & Aula 1 \\
3 & Rayleigh $+$ Beer: céu e poente & Aula 2 \\
4 & miragens: raios integrados em $n(z)$ & Aula 2 \\
\hline
\end{tabular}
\end{center}

Os Casos 1--2 são Monte Carlo de verdade com física de graduação — um
fecho computacional elegante para o curso.

\subsection*{Exercícios complementares e soluções}

Nas notas: T1--T5 e N1--N4. Soluções em \texttt{cap22\_solucoes.ipynb}
(\emph{não distribuir}): T1 e T2 com \texttt{sympy}; N4 (looming) é o
mais bonito. Pares: T1$\leftrightarrow$N1, T2$\leftrightarrow$N2,
T3$\leftrightarrow$N3, T4$\leftrightarrow$N4.

\subsection*{Conexões com o resto do curso}

Beer e massa de ar (Cap.~2); gotas e cristais (Caps.~6--7: o arco
exige chuva, o halo exige cirrus); halo como precursor frontal
(Cap.~12); inversões e miragens (Cap.~5); turbulência e cintilação
(Cap.~18); aerossol vulcânico e poentes (Cap.~21). O módulo extra de
Mudança Climática (Cap.~23) pode fechar o semestre após este.
